\documentclass[a0paper,portrait]{baposter}

\usepackage{algorithm}
\usepackage{algpseudocode}
\usepackage{amsmath}
\usepackage{amssymb}
\usepackage{booktabs}
% \usepackage{hyperref}
\usepackage{ifthen}
\usepackage{keyval}
\usepackage{mathtools}
\usepackage{placeins}
\usepackage{siunitx}
\usepackage{tabu}
\usepackage{tikz}
\usepackage{upgreek}

\usepackage{wrapfig}
\usepackage{lmodern}
\usepackage{tabularx,booktabs}
\newcolumntype{C}{>{\centering\arraybackslash}X} % centered version of "X" type
% \usepackage{url}
\usepackage{multirow}
\usepackage{array}

\usepackage[utf8]{inputenc} %unicode support
\usepackage[T1]{fontenc}

\usetikzlibrary{arrows.meta, backgrounds, calc, fit, positioning}

\DeclareMathOperator*{\binop}{\oplus}

\selectcolormodel{cmyk}

\graphicspath{{data/}}

\newcommand{\compresslist}{%
\setlength{\itemsep}{0pt}%
\setlength{\parskip}{0pt}%
\setlength{\parsep}{0pt}%
}

\newenvironment{boenumerate}
  {\begin{enumerate}\renewcommand\labelenumi{\textbf\theenumi.}}
  {\end{enumerate}}

% bold without changing font size
% https://tex.stackexchange.com/questions/23678/textbftext-without-increasing-the-length-of-the-text
\newsavebox\CBox
\def\textBF#1{\sbox\CBox{#1}\resizebox{\wd\CBox}{\ht\CBox}{\textbf{#1}}}


\begin{document}

\newcommand{\thetavec}{\ensuremath{\boldsymbol{\theta}}}
\newcommand{\fusop}{\ensuremath{\mathcal{F_{\thetavec}}}}
\newcommand{\mutanfeat}[1]{\ensuremath{\tilde{\mathbf{#1}}}}
\renewcommand{\m}{\ensuremath{\mathbf{m}}}
\newcommand{\n}{\ensuremath{\mathbf{n}}}
\newcommand{\q}{\ensuremath{\mathbf{q}}}
\renewcommand{\v}{\ensuremath{\mathbf{v}}}
\newcommand{\z}{\ensuremath{\mathbf{z}}}
\renewcommand{\A}{\ensuremath{A}}
\newcommand{\B}{\ensuremath{B}}
\newcommand{\CC}{\ensuremath{C}}
\newcommand{\R}{\ensuremath{\mathbb{R}}}
\newcommand{\T}{\ensuremath{\mathcal{T}}}
\newcommand{\Tfibre}{\ensuremath{\T_{i_1 \cdots i_{n - 1} \,:\, i_{n + 1} \cdots i_N}}}
\newcommand{\TsizeIn}[1]{\ensuremath{I_1 \times \cdots \times I_{n - 1} \times #1 \times I_{n + 1} \times \cdots \times I_N}}
\newcommand{\binopb}{\ensuremath{\sideset{}{_b}\binop}}
\newcommand{\binopseq}{\ensuremath{{(\binopb)}_{b = 1}^B}}
\newcommand{\binoppartition}{\ensuremath{{\{\tuckbranch\}}_b}}
\newcommand{\binoppartB}{\ensuremath{\mathbb{B}}}
\newcommand{\binoppartBseq}{\ensuremath{{(\binoppartB_b)}_{b = 1}^B}}
\newcommand{\tuckbranch}{\ensuremath{\T_r^{\q{}\v{}}}}
\newcommand{\hadamardqv}{\ensuremath{N_r\mutanfeat{q} \odot M_r\mutanfeat{v}}}
\newcommand{\qtWq}{\q^\intercal{} W_\q{}}
\newcommand{\vtWv}{\v^\intercal{} W_\v{}}

\definecolor{darkgreen}{cmyk}{0.8,0,0.8,0.45}
\definecolor{lightgreen}{cmyk}{0.8,0,0.8,0.25}
\definecolor{vector-magenta}{cmyk}{0,1,0,0}
\definecolor{vector-teal}{cmyk}{.67,0,.41,0}
\definecolor{vector-darkblue}{cmyk}{0.96,0.65,0.1,0.25}
\begin{poster}
{
grid=false,
headerborder=open, % Adds a border around the header of content boxes
colspacing=0.5em, % Column spacing
bgColorOne=white, % Background color for the gradient on the left side of the poster
bgColorTwo=white, % Background color for the gradient on the right side of the poster
borderColor=vector-magenta, % Border color
headerColorOne=vector-magenta, % Background color for the header in the content boxes (left side)
headerColorTwo=vector-magenta, % Background color for the header in the content boxes (right side)
headerFontColor=white, % Text color for the header text in the content boxes
boxColorOne=white, % Background color of the content boxes
textborder=rounded, %rectangle, % Format of the border around content boxes,
                    % can be: none, bars, coils, triangles, rectangle, rounded,
                    % roundedsmall, roundedright or faded
eyecatcher=false, % Set to false for ignoring the left logo in the title and move the title left
headerheight=0.1\textheight, % Height of the header
headershape=rounded, % Specify the rounded corner in the content box headers,
                     % can be: rectangle, small-rounded, roundedright,
                     % roundedleft or rounded
headershade=plain,
headerfont=\Large\textsf, % Large, bold and sans serif font in the headers of content boxes
%textfont={\setlength{\parindent}{1.5em}}, % Uncomment for paragraph indentation
linewidth=2pt % Width of the border lines around content boxes
}
{}
%
%----------------------------------------------------------------------------------------
%	TITLE AND AUTHOR NAME
%----------------------------------------------------------------------------------------
%
{
\textsf %Sans Serif
{Generalized Hadamard-Product\\
Fusion Operators for Visual Question Answering}
} % Poster title
{\sf\vspace{0em}\\
Brendan Duke\textsuperscript{1,2} and
Graham W.~Taylor\textsuperscript{1,2,3}
\vspace{0.1em}\\
\small{
\textsuperscript{1}University of Guelph,
\textsuperscript{2}Vector Institute, and 
\textsuperscript{3}Canadian Institute for Advanced Research
\vspace{0.2em}\\
\{bduke, gwtaylor\}@uoguelph.ca}
}
%{\includegraphics{logo}} % University/lab logo
{
\begin{minipage}[b]{0.15\linewidth}
\begin{minipage}[b]{\linewidth}
    \includegraphics[width=0.5\linewidth]{UofG_Cornerstone_cmyk.pdf}
\end{minipage}

\begin{minipage}[b]{\linewidth}
    \includegraphics[width=0.5\linewidth]{Vector_Logo_Bilingual_VERT_CMYK.pdf}
\end{minipage}
\end{minipage}
}

\headerbox{1.~Introduction}{name=introduction,column=0,row=0, span=3}{%
We propose a generalized class of multimodal fusion operators for the task of
visual question answering (VQA). We identify generalizations of existing
multimodal fusion operators based on the Hadamard product, and show that
specific non-trivial instantiations of this generalized fusion operator exhibit
superior performance in terms of OpenEnded accuracy on the VQA task. In
particular, we introduce Nonlinearity Ensembling, Feature Gating, and
post-fusion neural network layers as fusion operator components, culminating in
an absolute percentage point improvement of~$1.1\%$ on the VQA~2.0 test-dev set
over baseline fusion operators, which use the same features as input. We use
our findings as evidence that our generalized class of fusion operators could
lead to the discovery of even superior task-specific operators when used as a
search space in an architecture search over fusion operators.
}

\headerbox{2.~VQA Task}{name=vqa-task,column=1,below=introduction,span=2}{%
In general, we define a fusion operator for the VQA task as a function
$\fusop$, parametrized by~$\thetavec$, of the question feature vector~$\q$ and
the visual feature vector~$\v$. The fusion operator~$\fusop$ computes a vector
output, which is consumed downstream by a function~$g$, e.g.\ a linear layer
followed by a softmax layer, in order to model a probability distribution over
the answer~$y$ conditional on~$\q$ and~$\v$:

\begin{equation}
        p(y \mid \q \,, \v \,; \thetavec) = g(\fusop{(\q \,, \v)}) \,.\
\label{eqn:general-fusion-op}
\end{equation}

We looked for generalized expressions for the fusion operator~$\fusop$ based on
the Tucker Decomposition.
}

\headerbox{3.~Tucker Decomposition}{name=tucker-decomp,column=0,below=introduction,span=1}{%
\hspace{-65mm}\resizebox{1.9\linewidth}{!}{% NOTE(brendan): Source:
% https://tex.stackexchange.com/questions/29877/need-help-creating-a-3d-cube-from-a-2d-set-of-nodes-in-tikz

\makeatletter

% Standard Values for Parameters
\newcommand{\tikzcuboid@shiftx}{0}
\newcommand{\tikzcuboid@shifty}{0}
\newcommand{\tikzcuboid@dimx}{3}
\newcommand{\tikzcuboid@dimy}{3}
\newcommand{\tikzcuboid@dimz}{3}
\newcommand{\tikzcuboid@scale}{1}
\newcommand{\tikzcuboid@densityx}{1}
\newcommand{\tikzcuboid@densityy}{1}
\newcommand{\tikzcuboid@densityz}{1}
\newcommand{\tikzcuboid@rotation}{0}
\newcommand{\tikzcuboid@anglex}{0}
\newcommand{\tikzcuboid@angley}{90}
\newcommand{\tikzcuboid@anglez}{225}
\newcommand{\tikzcuboid@scalex}{1}
\newcommand{\tikzcuboid@scaley}{1}
\newcommand{\tikzcuboid@scalez}{sqrt(0.5)}
\newcommand{\tikzcuboid@linefront}{black}
\newcommand{\tikzcuboid@linetop}{black}
\newcommand{\tikzcuboid@lineright}{black}
\newcommand{\tikzcuboid@fillfront}{white}
\newcommand{\tikzcuboid@filltop}{white}
\newcommand{\tikzcuboid@fillright}{white}
\newcommand{\tikzcuboid@shaded}{N}
\newcommand{\tikzcuboid@shadecolor}{black}
\newcommand{\tikzcuboid@shadeperc}{25}
\newcommand{\tikzcuboid@emphedge}{N}
\newcommand{\tikzcuboid@emphstyle}{thick}
\newcommand{\foreachsteppingx}{%
\ifthenelse{\equal{\dimx}{1}}
    {\foreach \x in {\steppingx,...,\dimx}}
    {\foreach \x in {\steppingx,\secondx,...,\dimx}}
}
\newcommand{\foreachsteppingy}{%
\ifthenelse{\equal{\dimy}{1}}
    {\foreach \y in {\steppingy,...,\dimy}}
    {\foreach \y in {\steppingy,\secondy,...,\dimy}}
}
\newcommand{\foreachsteppingz}{%
\ifthenelse{\equal{\dimz}{1}}
    {\foreach \z in {\steppingz,...,\dimz}}
    {\foreach \z in {\steppingz,\secondz,...,\dimz}}
}

% Definition of Keys
\define@key{tikzcuboid}{shiftx}[\tikzcuboid@shiftx]{\renewcommand{\tikzcuboid@shiftx}{#1}}
\define@key{tikzcuboid}{shifty}[\tikzcuboid@shifty]{\renewcommand{\tikzcuboid@shifty}{#1}}
\define@key{tikzcuboid}{dimx}[\tikzcuboid@dimx]{\renewcommand{\tikzcuboid@dimx}{#1}}
\define@key{tikzcuboid}{dimy}[\tikzcuboid@dimy]{\renewcommand{\tikzcuboid@dimy}{#1}}
\define@key{tikzcuboid}{dimz}[\tikzcuboid@dimz]{\renewcommand{\tikzcuboid@dimz}{#1}}
\define@key{tikzcuboid}{scale}[\tikzcuboid@scale]{\renewcommand{\tikzcuboid@scale}{#1}}
\define@key{tikzcuboid}{densityx}[\tikzcuboid@densityx]{\renewcommand{\tikzcuboid@densityx}{#1}}
\define@key{tikzcuboid}{densityy}[\tikzcuboid@densityy]{\renewcommand{\tikzcuboid@densityy}{#1}}
\define@key{tikzcuboid}{densityz}[\tikzcuboid@densityz]{\renewcommand{\tikzcuboid@densityz}{#1}}
\define@key{tikzcuboid}{rotation}[\tikzcuboid@rotation]{\renewcommand{\tikzcuboid@rotation}{#1}}
\define@key{tikzcuboid}{anglex}[\tikzcuboid@anglex]{\renewcommand{\tikzcuboid@anglex}{#1}}
\define@key{tikzcuboid}{angley}[\tikzcuboid@angley]{\renewcommand{\tikzcuboid@angley}{#1}}
\define@key{tikzcuboid}{anglez}[\tikzcuboid@anglez]{\renewcommand{\tikzcuboid@anglez}{#1}}
\define@key{tikzcuboid}{scalex}[\tikzcuboid@scalex]{\renewcommand{\tikzcuboid@scalex}{#1}}
\define@key{tikzcuboid}{scaley}[\tikzcuboid@scaley]{\renewcommand{\tikzcuboid@scaley}{#1}}
\define@key{tikzcuboid}{scalez}[\tikzcuboid@scalez]{\renewcommand{\tikzcuboid@scalez}{#1}}
\define@key{tikzcuboid}{linefront}[\tikzcuboid@linefront]{\renewcommand{\tikzcuboid@linefront}{#1}}
\define@key{tikzcuboid}{linetop}[\tikzcuboid@linetop]{\renewcommand{\tikzcuboid@linetop}{#1}}
\define@key{tikzcuboid}{lineright}[\tikzcuboid@lineright]{\renewcommand{\tikzcuboid@lineright}{#1}}
\define@key{tikzcuboid}{fillfront}[\tikzcuboid@fillfront]{\renewcommand{\tikzcuboid@fillfront}{#1}}
\define@key{tikzcuboid}{filltop}[\tikzcuboid@filltop]{\renewcommand{\tikzcuboid@filltop}{#1}}
\define@key{tikzcuboid}{fillright}[\tikzcuboid@fillright]{\renewcommand{\tikzcuboid@fillright}{#1}}
\define@key{tikzcuboid}{shaded}[\tikzcuboid@shaded]{\renewcommand{\tikzcuboid@shaded}{#1}}
\define@key{tikzcuboid}{shadecolor}[\tikzcuboid@shadecolor]{\renewcommand{\tikzcuboid@shadecolor}{#1}}
\define@key{tikzcuboid}{shadeperc}[\tikzcuboid@shadeperc]{\renewcommand{\tikzcuboid@shadeperc}{#1}}
\define@key{tikzcuboid}{emphedge}[\tikzcuboid@emphedge]{\renewcommand{\tikzcuboid@emphedge}{#1}}
\define@key{tikzcuboid}{emphstyle}[\tikzcuboid@emphstyle]{\renewcommand{\tikzcuboid@emphstyle}{#1}}
% Commands
\newcommand{\tikzcuboid}[1]{
    \setkeys{tikzcuboid}{#1} % Process Keys passed to command
    \pgfmathsetmacro{\vectorxx}{\tikzcuboid@scalex*cos(\tikzcuboid@anglex)}
    \pgfmathsetmacro{\vectorxy}{\tikzcuboid@scalex*sin(\tikzcuboid@anglex)}
    \pgfmathsetmacro{\vectoryx}{\tikzcuboid@scaley*cos(\tikzcuboid@angley)}
    \pgfmathsetmacro{\vectoryy}{\tikzcuboid@scaley*sin(\tikzcuboid@angley)}
    \pgfmathsetmacro{\vectorzx}{\tikzcuboid@scalez*cos(\tikzcuboid@anglez)}
    \pgfmathsetmacro{\vectorzy}{\tikzcuboid@scalez*sin(\tikzcuboid@anglez)}
    \begin{scope}[xshift=\tikzcuboid@shiftx, yshift=\tikzcuboid@shifty, scale=\tikzcuboid@scale, rotate=\tikzcuboid@rotation, x={(\vectorxx,\vectorxy)}, y={(\vectoryx,\vectoryy)}, z={(\vectorzx,\vectorzy)}]
    \pgfmathsetmacro{\steppingx}{1/\tikzcuboid@densityx}
    \pgfmathsetmacro{\steppingy}{1/\tikzcuboid@densityy}
    \pgfmathsetmacro{\steppingz}{1/\tikzcuboid@densityz}
    \newcommand{\dimx}{\tikzcuboid@dimx}
    \newcommand{\dimy}{\tikzcuboid@dimy}
    \newcommand{\dimz}{\tikzcuboid@dimz}
    \pgfmathsetmacro{\secondx}{2*\steppingx}
    \pgfmathsetmacro{\secondy}{2*\steppingy}
    \pgfmathsetmacro{\secondz}{2*\steppingz}
    \foreachsteppingx
    {
        \foreachsteppingy
        {   \pgfmathsetmacro{\lowx}{(\x-\steppingx)}
            \pgfmathsetmacro{\lowy}{(\y-\steppingy)}
            \filldraw[fill=\tikzcuboid@fillfront,draw=\tikzcuboid@linefront] (\lowx,\lowy,\dimz) -- (\lowx,\y,\dimz) -- (\x,\y,\dimz) -- (\x,\lowy,\dimz) -- cycle;

        }
    }
    \foreachsteppingx
    {
        \foreachsteppingz
        {   \pgfmathsetmacro{\lowx}{(\x-\steppingx)}
            \pgfmathsetmacro{\lowz}{(\z-\steppingz)}
            \filldraw[fill=\tikzcuboid@filltop,draw=\tikzcuboid@linetop] (\lowx,\dimy,\lowz) -- (\lowx,\dimy,\z) -- (\x,\dimy,\z) -- (\x,\dimy,\lowz) -- cycle;
        }
    }
    \foreachsteppingy
    {
        \foreachsteppingz
        {   \pgfmathsetmacro{\lowy}{(\y-\steppingy)}
            \pgfmathsetmacro{\lowz}{(\z-\steppingz)}
            \filldraw[fill=\tikzcuboid@fillright,draw=\tikzcuboid@lineright] (\dimx,\lowy,\lowz) -- (\dimx,\lowy,\z) -- (\dimx,\y,\z) -- (\dimx,\y,\lowz) -- cycle;
        }
    }
    \ifthenelse{\equal{\tikzcuboid@emphedge}{Y}}%
        {\draw[\tikzcuboid@emphstyle](0,\dimy,0) -- (\dimx,\dimy,0) -- (\dimx,\dimy,\dimz) -- (0,\dimy,\dimz) -- cycle;%
        \draw[\tikzcuboid@emphstyle] (0,0,\dimz) -- (0,\dimy,\dimz) -- (\dimx,\dimy,\dimz) -- (\dimx,0,\dimz) -- cycle;%
        \draw[\tikzcuboid@emphstyle](\dimx,0,0) -- (\dimx,\dimy,0) -- (\dimx,\dimy,\dimz) -- (\dimx,0,\dimz) -- cycle;%
        }%
        {}
    \end{scope}
}

\newcommand{\qvectikzdim}{5}
\newcommand{\vvectikzdim}{4}
\newcommand{\cuboidscale}{0.35}


\begin{tikzpicture}
        % NOTE(brendan): Tc
        \tikzcuboid{%
            shiftx=0cm,%
            shifty=0cm,%
            scale=\cuboidscale,%
            rotation=0,%
            densityx=1,%
            densityy=1,%
            densityz=1,%
            dimx=\qvectikzdim,%
            dimy=\vvectikzdim,%
            dimz=3,%
            linefront=black,%
            linetop=black,%
            lineright=black,%
            fillfront=red!50,%
            filltop=green!50,%
            fillright=blue!50,%
            anglex=180,%
            angley=90,%
            anglez=225,%
            scalex=1,%
            scaley=0.8,%
            scalez=0.6,%
            emphedge=N,%
        }

        % NOTE(brendan):
        \tikzcuboid{%
            shiftx=2.0cm,%
            shifty=0cm,%
            scale=\cuboidscale,%
            rotation=0,%
            densityx=1,%
            densityy=1,%
            densityz=1,%
            dimx=1,%
            dimy=\vvectikzdim,%
            dimz=3,%
            linefront=black,%
            linetop=black,%
            lineright=black,%
            fillfront=red!50,%
            filltop=green!50,%
            fillright=blue!50,%
            anglex=180,%
            angley=90,%
            anglez=225,%
            scalex=1,%
            scaley=0.8,%
            scalez=0.6,%
            emphedge=N,%
        }

        \tikzcuboid{%
            shiftx=10pt,%
            shifty=-15mm,%
            scale=\cuboidscale,%
            rotation=0,%
            densityx=1,%
            densityy=1,%
            densityz=1,%
            dimx=\qvectikzdim,%
            dimy=1,%
            dimz=1,%
            linefront=black,%
            linetop=black,%
            lineright=black,%
            fillfront=red!50,%
            filltop=red!50,%
            fillright=red!50,%
            anglex=180,%
            angley=90,%
            anglez=225,%
            scalex=1,%
            scaley=0.8,%
            scalez=0.6,%
            emphedge=N,%
        }

        \newcommand{\qbracesize}{8.5ex}
        \newcommand{\vbracesize}{5ex}

        \coordinate (qbrace) at (-0.35, -1.77);
        \node at (qbrace) {\makebox[\qbracesize]{\upbracefill}};
        \node at (qbrace) [below] {$t_q$};

        \coordinate (Tcq brace) at (-0.42, -0.18);
        \node at (Tcq brace) {\rotatebox{180}{\makebox[\qbracesize]{\upbracefill}}};
        \node[inner sep=1pt] at (Tcq brace) [above, fill=white, fill opacity=0.8, text opacity=1, yshift=1.1mm] {$t_q$};

        \node (Tc label) at (Tcq brace) [xshift=-3mm, yshift=18mm] {$\T_c$};
        \node (Tcq label) at (Tc label) [xshift=27mm] {$\T_c^\q$};
        \node (vtWv label) at (Tc label) [xshift=45mm] {$\vtWv$};
        \node[inner sep=0pt] (qtWq label) at (Tc label) [yshift=-43mm, xshift=3mm] {$\qtWq$};
        \node (Tcqv label) at (Tc label) [xshift=65mm] {$\T_c^{\q\v}$};

        \newcommand{\Tcvy}{0.56}
        \coordinate (Tcv brace) at (-1.9, \Tcvy);
        \node at (Tcv brace) {\rotatebox{270}{\makebox[\vbracesize]{\upbracefill}}};
        \node at (Tcv brace) [left] {$t_v$};

        \coordinate (Tco brace) at (0.33, 1.05);
        \node at (Tco brace) {\rotatebox{140}{\makebox[2ex]{\upbracefill}}};
        \node at (Tco brace) [right, yshift=2.0mm, xshift=-1.0mm] {$t_o$};

        \node at (qbrace) [yshift=10mm] {$\times_1$};

        \tikzcuboid{%
            shiftx=3.8cm,%
            shifty=-7pt,%
            scale=\cuboidscale,%
            rotation=0,%
            densityx=1,%
            densityy=1,%
            densityz=1,%
            dimx=1,%
            dimy=\vvectikzdim,%
            dimz=1,%
            linefront=black,%
            linetop=black,%
            lineright=black,%
            fillfront=blue!50,%
            filltop=blue!50,%
            fillright=blue!50,%
            anglex=180,%
            angley=90,%
            anglez=225,%
            scalex=1,%
            scaley=0.8,%
            scalez=0.6,%
            emphedge=N,%
        }

        \coordinate (vbrace) at (4.1, 0.18);
        \node at (vbrace) {\rotatebox{90}{\makebox[\vbracesize]{\upbracefill}}};
        \node at (vbrace) [right] {$t_v$};

        \draw[->, thick, >= stealth] (0.85, 0.25) -- (1.35, 0.25) {};

        \node at (vbrace) [xshift=-11.5mm, yshift=0mm] {$\times_2$};

        \coordinate (tcqv pos) at (5.7, 0.31);

        \draw[->, thick, >= stealth] (4.63, 0.25) -- (5.13, 0.25) {};

        \tikzcuboid{%
            shiftx=5.76cm,%
            shifty=0.3cm,%
            scale=\cuboidscale,%
            rotation=0,%
            densityx=1,%
            densityy=1,%
            densityz=1,%
            dimx=1,%
            dimy=1,%
            dimz=3,%
            linefront=black,%
            linetop=black,%
            lineright=black,%
            fillfront=red!50,%
            filltop=green!50,%
            fillright=blue!50,%
            anglex=180,%
            angley=90,%
            anglez=225,%
            scalex=1,%
            scaley=0.8,%
            scalez=0.6,%
            emphedge=N,%
        }
\end{tikzpicture}
}
A geometric representation of the first two $n$-mode products of the
right hand side Tucker decomposition
$\T \approx \T_c \times_1 \A \times_2 \B \times_3 \CC$,
where~$A = \qtWq$ and~$B = \vtWv$, i.e.\ $\T_c \times_1 \qtWq \times_2 \vtWv$.
From left to right, the vector~$\qtWq$ is first combined with the core
tensor~$\T_c$ by taking the inner product along the dimension of
size~$t_q$ between~$\qtWq$ and each of the~$t_v \times t_o$ fibers
composing~$\T_c$, producing the matrix $\T_c^\q$.  Then,~$\vtWv$ is
combined with~$\T_c^\q$ by taking the inner product between~$\vtWv$
and the~$t_o$ columns of~$\T_c^\q$, producing the
vector~$\T_c^{\q{}\v{}}$.
}

\headerbox{4.~Comparison with MUTAN}{name=comparison-mutan,column=1,below=vqa-task,span=2}{%
A comparison between the MUTAN fusion operator of~\cite{ben2017mutan} and our
generalized fusion operator.

Both MUTAN and our fusion operator take the features~$M_r\mutanfeat{q}$
and~$N_r\mutanfeat{v}$ as input. MUTAN approximates the Tucker decomposition by
constraining the bilinear interaction between the vectors~$\mutanfeat{q}$
and~$\mutanfeat{v}$ to be of rank~$R$.

\begin{center}
% NOTE(brendan): Source:
% https://tex.stackexchange.com/questions/29877/need-help-creating-a-3d-cube-from-a-2d-set-of-nodes-in-tikz

\makeatletter

% Standard Values for Parameters
\newcommand{\tikzcuboid@shiftx}{0}
\newcommand{\tikzcuboid@shifty}{0}
\newcommand{\tikzcuboid@dimx}{3}
\newcommand{\tikzcuboid@dimy}{3}
\newcommand{\tikzcuboid@dimz}{3}
\newcommand{\tikzcuboid@scale}{1}
\newcommand{\tikzcuboid@densityx}{1}
\newcommand{\tikzcuboid@densityy}{1}
\newcommand{\tikzcuboid@densityz}{1}
\newcommand{\tikzcuboid@rotation}{0}
\newcommand{\tikzcuboid@anglex}{0}
\newcommand{\tikzcuboid@angley}{90}
\newcommand{\tikzcuboid@anglez}{225}
\newcommand{\tikzcuboid@scalex}{1}
\newcommand{\tikzcuboid@scaley}{1}
\newcommand{\tikzcuboid@scalez}{sqrt(0.5)}
\newcommand{\tikzcuboid@linefront}{black}
\newcommand{\tikzcuboid@linetop}{black}
\newcommand{\tikzcuboid@lineright}{black}
\newcommand{\tikzcuboid@fillfront}{white}
\newcommand{\tikzcuboid@filltop}{white}
\newcommand{\tikzcuboid@fillright}{white}
\newcommand{\tikzcuboid@shaded}{N}
\newcommand{\tikzcuboid@shadecolor}{black}
\newcommand{\tikzcuboid@shadeperc}{25}
\newcommand{\tikzcuboid@emphedge}{N}
\newcommand{\tikzcuboid@emphstyle}{thick}
\newcommand{\foreachsteppingx}{%
\ifthenelse{\equal{\dimx}{1}}
    {\foreach \x in {\steppingx,...,\dimx}}
    {\foreach \x in {\steppingx,\secondx,...,\dimx}}
}
\newcommand{\foreachsteppingy}{%
\ifthenelse{\equal{\dimy}{1}}
    {\foreach \y in {\steppingy,...,\dimy}}
    {\foreach \y in {\steppingy,\secondy,...,\dimy}}
}
\newcommand{\foreachsteppingz}{%
\ifthenelse{\equal{\dimz}{1}}
    {\foreach \z in {\steppingz,...,\dimz}}
    {\foreach \z in {\steppingz,\secondz,...,\dimz}}
}

% Definition of Keys
\define@key{tikzcuboid}{shiftx}[\tikzcuboid@shiftx]{\renewcommand{\tikzcuboid@shiftx}{#1}}
\define@key{tikzcuboid}{shifty}[\tikzcuboid@shifty]{\renewcommand{\tikzcuboid@shifty}{#1}}
\define@key{tikzcuboid}{dimx}[\tikzcuboid@dimx]{\renewcommand{\tikzcuboid@dimx}{#1}}
\define@key{tikzcuboid}{dimy}[\tikzcuboid@dimy]{\renewcommand{\tikzcuboid@dimy}{#1}}
\define@key{tikzcuboid}{dimz}[\tikzcuboid@dimz]{\renewcommand{\tikzcuboid@dimz}{#1}}
\define@key{tikzcuboid}{scale}[\tikzcuboid@scale]{\renewcommand{\tikzcuboid@scale}{#1}}
\define@key{tikzcuboid}{densityx}[\tikzcuboid@densityx]{\renewcommand{\tikzcuboid@densityx}{#1}}
\define@key{tikzcuboid}{densityy}[\tikzcuboid@densityy]{\renewcommand{\tikzcuboid@densityy}{#1}}
\define@key{tikzcuboid}{densityz}[\tikzcuboid@densityz]{\renewcommand{\tikzcuboid@densityz}{#1}}
\define@key{tikzcuboid}{rotation}[\tikzcuboid@rotation]{\renewcommand{\tikzcuboid@rotation}{#1}}
\define@key{tikzcuboid}{anglex}[\tikzcuboid@anglex]{\renewcommand{\tikzcuboid@anglex}{#1}}
\define@key{tikzcuboid}{angley}[\tikzcuboid@angley]{\renewcommand{\tikzcuboid@angley}{#1}}
\define@key{tikzcuboid}{anglez}[\tikzcuboid@anglez]{\renewcommand{\tikzcuboid@anglez}{#1}}
\define@key{tikzcuboid}{scalex}[\tikzcuboid@scalex]{\renewcommand{\tikzcuboid@scalex}{#1}}
\define@key{tikzcuboid}{scaley}[\tikzcuboid@scaley]{\renewcommand{\tikzcuboid@scaley}{#1}}
\define@key{tikzcuboid}{scalez}[\tikzcuboid@scalez]{\renewcommand{\tikzcuboid@scalez}{#1}}
\define@key{tikzcuboid}{linefront}[\tikzcuboid@linefront]{\renewcommand{\tikzcuboid@linefront}{#1}}
\define@key{tikzcuboid}{linetop}[\tikzcuboid@linetop]{\renewcommand{\tikzcuboid@linetop}{#1}}
\define@key{tikzcuboid}{lineright}[\tikzcuboid@lineright]{\renewcommand{\tikzcuboid@lineright}{#1}}
\define@key{tikzcuboid}{fillfront}[\tikzcuboid@fillfront]{\renewcommand{\tikzcuboid@fillfront}{#1}}
\define@key{tikzcuboid}{filltop}[\tikzcuboid@filltop]{\renewcommand{\tikzcuboid@filltop}{#1}}
\define@key{tikzcuboid}{fillright}[\tikzcuboid@fillright]{\renewcommand{\tikzcuboid@fillright}{#1}}
\define@key{tikzcuboid}{shaded}[\tikzcuboid@shaded]{\renewcommand{\tikzcuboid@shaded}{#1}}
\define@key{tikzcuboid}{shadecolor}[\tikzcuboid@shadecolor]{\renewcommand{\tikzcuboid@shadecolor}{#1}}
\define@key{tikzcuboid}{shadeperc}[\tikzcuboid@shadeperc]{\renewcommand{\tikzcuboid@shadeperc}{#1}}
\define@key{tikzcuboid}{emphedge}[\tikzcuboid@emphedge]{\renewcommand{\tikzcuboid@emphedge}{#1}}
\define@key{tikzcuboid}{emphstyle}[\tikzcuboid@emphstyle]{\renewcommand{\tikzcuboid@emphstyle}{#1}}
% Commands
\newcommand{\tikzcuboid}[1]{
    \setkeys{tikzcuboid}{#1} % Process Keys passed to command
    \pgfmathsetmacro{\vectorxx}{\tikzcuboid@scalex*cos(\tikzcuboid@anglex)}
    \pgfmathsetmacro{\vectorxy}{\tikzcuboid@scalex*sin(\tikzcuboid@anglex)}
    \pgfmathsetmacro{\vectoryx}{\tikzcuboid@scaley*cos(\tikzcuboid@angley)}
    \pgfmathsetmacro{\vectoryy}{\tikzcuboid@scaley*sin(\tikzcuboid@angley)}
    \pgfmathsetmacro{\vectorzx}{\tikzcuboid@scalez*cos(\tikzcuboid@anglez)}
    \pgfmathsetmacro{\vectorzy}{\tikzcuboid@scalez*sin(\tikzcuboid@anglez)}
    \begin{scope}[xshift=\tikzcuboid@shiftx, yshift=\tikzcuboid@shifty, scale=\tikzcuboid@scale, rotate=\tikzcuboid@rotation, x={(\vectorxx,\vectorxy)}, y={(\vectoryx,\vectoryy)}, z={(\vectorzx,\vectorzy)}]
    \pgfmathsetmacro{\steppingx}{1/\tikzcuboid@densityx}
    \pgfmathsetmacro{\steppingy}{1/\tikzcuboid@densityy}
    \pgfmathsetmacro{\steppingz}{1/\tikzcuboid@densityz}
    \newcommand{\dimx}{\tikzcuboid@dimx}
    \newcommand{\dimy}{\tikzcuboid@dimy}
    \newcommand{\dimz}{\tikzcuboid@dimz}
    \pgfmathsetmacro{\secondx}{2*\steppingx}
    \pgfmathsetmacro{\secondy}{2*\steppingy}
    \pgfmathsetmacro{\secondz}{2*\steppingz}
    \foreachsteppingx
    {
        \foreachsteppingy
        {   \pgfmathsetmacro{\lowx}{(\x-\steppingx)}
            \pgfmathsetmacro{\lowy}{(\y-\steppingy)}
            \filldraw[fill=\tikzcuboid@fillfront,draw=\tikzcuboid@linefront] (\lowx,\lowy,\dimz) -- (\lowx,\y,\dimz) -- (\x,\y,\dimz) -- (\x,\lowy,\dimz) -- cycle;

        }
    }
    \foreachsteppingx
    {
        \foreachsteppingz
        {   \pgfmathsetmacro{\lowx}{(\x-\steppingx)}
            \pgfmathsetmacro{\lowz}{(\z-\steppingz)}
            \filldraw[fill=\tikzcuboid@filltop,draw=\tikzcuboid@linetop] (\lowx,\dimy,\lowz) -- (\lowx,\dimy,\z) -- (\x,\dimy,\z) -- (\x,\dimy,\lowz) -- cycle;
        }
    }
    \foreachsteppingy
    {
        \foreachsteppingz
        {   \pgfmathsetmacro{\lowy}{(\y-\steppingy)}
            \pgfmathsetmacro{\lowz}{(\z-\steppingz)}
            \filldraw[fill=\tikzcuboid@fillright,draw=\tikzcuboid@lineright] (\dimx,\lowy,\lowz) -- (\dimx,\lowy,\z) -- (\dimx,\y,\z) -- (\dimx,\y,\lowz) -- cycle;
        }
    }
    \ifthenelse{\equal{\tikzcuboid@emphedge}{Y}}%
        {\draw[\tikzcuboid@emphstyle](0,\dimy,0) -- (\dimx,\dimy,0) -- (\dimx,\dimy,\dimz) -- (0,\dimy,\dimz) -- cycle;%
        \draw[\tikzcuboid@emphstyle] (0,0,\dimz) -- (0,\dimy,\dimz) -- (\dimx,\dimy,\dimz) -- (\dimx,0,\dimz) -- cycle;%
        \draw[\tikzcuboid@emphstyle](\dimx,0,0) -- (\dimx,\dimy,0) -- (\dimx,\dimy,\dimz) -- (\dimx,0,\dimz) -- cycle;%
        }%
        {}
    \end{scope}
}

\newcommand{\qvectikzdim}{5}
\newcommand{\vvectikzdim}{4}
\newcommand{\cuboidscale}{0.35}


\newcommand{\outveccuboid}[2]{%
        \tikzcuboid{%
            shiftx=#1,%
            shifty=#2,%
            scale=\cuboidscale,%
            rotation=0,%
            densityx=1,%
            densityy=1,%
            densityz=1,%
            dimx=1,%
            dimy=1,%
            dimz=3,%
            linefront=black,%
            linetop=black,%
            lineright=black,%
            fillfront=red!50,%
            filltop=green!50,%
            fillright=blue!50,%
            anglex=180,%
            angley=90,%
            anglez=225,%
            scalex=1,%
            scaley=0.8,%
            scalez=0.6,%
            emphedge=N,%
        }
}

\begin{tikzpicture}
        [hadamard/.style={circle,
                          inner sep=0pt,
                          minimum size=6mm,
                          thick,
                          draw=blue!75,
                          fill=blue!20},
         oplussty/.style={circle,
                          inner sep=0pt,
                          minimum size=6mm,
                          thick,
                          draw=cyan!80,
                          fill=cyan!50},
         pre/.style={<-, shorten <= 1pt, >=stealth, semithick},
         post/.style={->, shorten >= 1pt, >=stealth, semithick},
         sumsty/.style={circle,
                        inner sep=0pt,
                        minimum size=6mm,
                        thick,
                        draw=orange!80,
                        fill=orange!50}]

        % Nr
        \newcommand{\nposx}{0.5mm}
        \newcommand{\nposy}{0cm}

        \newcommand{\mposx}{2.15cm}
        \newcommand{\mposy}{-1.0cm}

        \coordinate (mcoord) at (\mposx, \mposy);
        \coordinate (ncoord) at (\nposx, \nposy);

        \coordinate (mtopleft) at (\mposx -1.9cm, \mposy + 0.25cm);
        \coordinate (mbottomright) at (\mposx + 0.5cm, \mposy - 1.5cm);
        \coordinate (ntopleft) at (\nposx -2.0cm, \nposy + 1.15cm);
        \coordinate (nbottomright) at (\nposx + 0.5cm, \nposy - 0.4cm);

        \newcommand{\qtopcolour}{green}
        \newcommand{\qcolourright}{blue}
        \newcommand{\vtopcolour}{green}
        \newcommand{\vcolourfront}{red}
        \tikzcuboid{%
            shiftx=\nposx,%
            shifty=\nposy,%
            scale=\cuboidscale,%
            rotation=0,%
            densityx=1,%
            densityy=1,%
            densityz=1,%
            dimx=1,%
            dimy=\vvectikzdim,%
            dimz=3,%
            linefront=black,%
            linetop=black,%
            lineright=black,%
            fillfront=\vcolourfront!50,%
            filltop=\vtopcolour!50,%
            fillright=blue!50,%
            anglex=180,%
            angley=90,%
            anglez=225,%
            scalex=1,%
            scaley=0.8,%
            scalez=0.6,%
            emphedge=N,%
        }

        \tikzcuboid{%
            shiftx=\mposx,%
            shifty=\mposy,%
            scale=\cuboidscale,%
            rotation=0,%
            densityx=1,%
            densityy=1,%
            densityz=1,%
            dimx=\qvectikzdim,%
            dimy=1,%
            dimz=3,%
            linefront=black,%
            linetop=black,%
            lineright=black,%
            fillfront=red!50,%
            filltop=\qtopcolour!50,%
            fillright=\qcolourright!50,%
            anglex=180,%
            angley=90,%
            anglez=225,%
            scalex=1,%
            scaley=0.8,%
            scalez=0.6,%
            emphedge=N,%
        }

        \newcommand{\vposx}{\nposx - 1.25cm}
        \newcommand{\vposy}{\nposy - 0.24cm}
        \tikzcuboid{%
            shiftx=\vposx,%
            shifty=\vposy,%
            scale=\cuboidscale,%
            rotation=0,%
            densityx=1,%
            densityy=1,%
            densityz=1,%
            dimx=1,%
            dimy=\vvectikzdim,%
            dimz=1,%
            linefront=black,%
            linetop=black,%
            lineright=black,%
            fillfront=blue!50,%
            filltop=blue!50,%
            fillright=blue!50,%
            anglex=180,%
            angley=90,%
            anglez=225,%
            scalex=1,%
            scaley=0.8,%
            scalez=0.6,%
            emphedge=N,%
        }

        \newcommand{\qposx}{\mposx + 0.35cm}
        \newcommand{\qposy}{\mposy - 1.3cm}
        \tikzcuboid{%
            shiftx=\qposx,%
            shifty=\qposy,%
            scale=\cuboidscale,%
            rotation=0,%
            densityx=1,%
            densityy=1,%
            densityz=1,%
            dimx=\qvectikzdim,%
            dimy=1,%
            dimz=1,%
            linefront=black,%
            linetop=black,%
            lineright=black,%
            fillfront=red!50,%
            filltop=red!50,%
            fillright=red!50,%
            anglex=180,%
            angley=90,%
            anglez=225,%
            scalex=1,%
            scaley=0.8,%
            scalez=0.6,%
            emphedge=N,%
        }

        \node (times1) at (mcoord) [xshift=-4mm, yshift=-7mm] {$\times_1$};
        \node (times2) at (ncoord) [xshift=-7mm, yshift=2.5mm] {$\times_2$};

        \node[hadamard, xshift=-0.5cm, yshift=1.25cm] (mutanhadamard) at (mcoord) {$\odot$}
                edge [pre] ($(mtopleft) + (14.0mm, 0.2cm)$)
                edge [pre] ($(nbottomright) + (0.2cm, 6.5mm)$);

        \node (mrqhadamardnrv) at (mutanhadamard) [xshift=4mm, yshift=7mm] {$M_r\mutanfeat{q} \odot N_r\mutanfeat{v}$}
                edge [pre, shorten <= -2pt] (mutanhadamard);

        \node (Rmany) [right=of mutanhadamard, xshift=0mm, yshift=-9mm] {\rotatebox{90}{\makebox[20ex]{\upbracefill}}};

        \newcommand{\Rmanytextwidth}{15mm}
        \newcommand{\Rmanydescrip}{Repeat $R$ times}
        \node (Rmanytext) [right=of Rmany, xshift=-11mm, text width=\Rmanytextwidth] {\Rmanydescrip};

        \node[matrix] (mutanR) [right=of Rmany, xshift=-8mm]
                {%
                        \node (r1) {};\\[2mm]
                        \node (r2) [xshift=6mm] {};\\[2mm]
                        \node (r3) {};\\
                };

        \node[sumsty] (sum) [right=of mutanR, xshift=-0.5cm] {$\Sigma$}
              edge [pre, out=90, in=90] (r1) {}
              edge [pre, out=180, in=0] (r2) {}
              edge [pre, out=270, in=270] (r3) {};

        \node (tcqv) [right=of sum, xshift=-0.5cm] {} edge [pre] (sum) {};

        \pgfgetlastxy{\XCoord}{\YCoord};
        \outveccuboid{\XCoord + 5mm}{\YCoord}

        \begin{scope}[on background layer]
                \node (Mbackground) [fill=red!10, rounded corners, fit=(mtopleft) (mbottomright)] {};
        \end{scope}

        \begin{scope}[on background layer]
                \node (Nbackground) [fill=blue!10, rounded corners, fit=(ntopleft) (nbottomright)] {};
        \end{scope}

        \node[inner sep=0pt] (qlabel) at (mtopleft) [yshift=-1.5cm, xshift=2mm] {$\mutanfeat{q}$};
        \node[inner sep=1pt] (Mlabel) at (mtopleft)
                [yshift=-2mm, xshift=13mm, fill=white, fill opacity=0.9, text opacity=1] {$M_r$};

        \node[inner sep=0pt] (vlabel) at (ntopleft) [yshift=-0.2cm, xshift=0.15cm] {$\mutanfeat{v}$};
        \node[inner sep=1pt] (Nlabel) at (ntopleft)
                [yshift=-8.5mm, xshift=20mm, fill=white, fill opacity=0.9, text opacity=1] {$N_r$};

        \node[matrix] (legend) at (current bounding box.east) [xshift=3.5cm] {%
                \node[hadamard, label=right:Element-wise multiplication.] {$\odot$};\\[2mm]
                \node[sumsty, label=right: Element-wise sum.] {$\Sigma$};\\[2mm]
                \node[oplussty, label=right: Binary operator~$b$.] {$\binopb$};\\[2mm]
                \node[minimum size=6mm, label={[xshift=1mm]right: $\qtWq$}] {$\mutanfeat{q}$};\\[2mm]
                \node[minimum size=6mm, label={[xshift=1mm]right: $\vtWv$}] {$\mutanfeat{v}$};\\
        };

        \begin{scope}[on background layer]
                \node (r1) [fill=black!10, rounded corners, fit=(legend)] {};
        \end{scope}


        \newcommand{\ourstopleftx}{-1.5cm}
        \newcommand{\ourstoplefty}{-2.5cm}
        \coordinate (ours top left) at (\ourstopleftx, \ourstoplefty);

        \node (q1) [below=of ours top left] {$M_{r}\mutanfeat{q}$};
        \node (v1) [below=of q1] {$N_{r}\mutanfeat{v}$};
        \node[hadamard, xshift=1.5cm] (h1) at ($(q1)!0.5!(v1)$) {$\odot$}
                 edge [pre] node [below] {$f_{\mathrm{rv}}$} (v1) {}
                 edge [pre] node [above, xshift=1mm] {$f_{\mathrm{rq}}$} (q1) {};

        \begin{scope}
                [cnode/.style={draw=black,fill=#1,minimum width=3mm,circle},
                 scale=0.5,
                 shift={($(h1) + (2.0cm, 3.0cm)$)}]
                \newcommand{\layertwoxshift}{2}
                \node at (0,-3.75) {$\vdots$};
                \node at (\layertwoxshift,-3.75) {$\vdots$};
                \foreach \x in {1,...,4}
                {
                        \pgfmathparse{\x<4 ? \x : "n"}
                        \node[cnode=gray!50] (x-\x) at (0,{-\x-div(\x,4)}) {};
                        \node[cnode=gray!50] (p-\x) at (\layertwoxshift,{-\x-div(\x,4)}) {};
                }
                \foreach \x in {1,...,4}
                {   \foreach \y in {1,...,4}
                {   \draw (x-\x) -- (p-\y);
                }
                }

                \draw[post] (h1) -- (x-3);
        \end{scope}

        \node (oursRmany) [right=of p-1, xshift=-9mm, yshift=-10mm] {\rotatebox{90}{\makebox[12ex]{\upbracefill}}};
        \node (oursRmanytext) [right=of oursRmany, xshift=-11mm, text width=\Rmanytextwidth] {\Rmanydescrip};

        \pgfgetlastxy{\XCoord}{\YCoord};
        \foreach \y/\ytext in {0cm/0, -0.75cm/1, -1.5cm/2, -2.25cm/3}
        {
                \node (outvec\ytext) at (\XCoord + 30mm, \YCoord + 11mm + \y) {}
                        edge [pre, out=180, in=0, shorten <= 8pt, shorten >= -10pt] (oursRmanytext) {};
                \outveccuboid{\XCoord + 30mm}{\YCoord + 10mm + \y}
        }

        \coordinate (outvec01mid) at ($(outvec0)!0.5!(outvec1)$);
        \newcommand{\oplusedgexshift}{3pt};
        \newcommand{\oplusinangleLO}{145}
        \newcommand{\oplusinangleHI}{215}
        \node[oplussty, label={[xshift=0.5mm, yshift=0mm]center:$\binop_1$}]
                (oplus1) [right=of outvec01mid] {}
                edge [pre, out=\oplusinangleLO, in=0, shorten >= \oplusedgexshift] (outvec0) {}
                edge [pre, out=\oplusinangleHI, in=0, shorten >= \oplusedgexshift] (outvec1) {};

        \coordinate (outvec23mid) at ($(outvec2)!0.5!(outvec3)$);
        \node[oplussty, label={[xshift=0.5mm, yshift=0mm]center:$\binop_2$}]
                (oplus2) [right=of outvec23mid] {}
                edge [pre, out=\oplusinangleLO, in=0, shorten >= \oplusedgexshift] (outvec2) {}
                edge [pre, out=\oplusinangleHI, in=0, shorten >= \oplusedgexshift] (outvec3) {};

        \coordinate (oplusmid) at ($(oplus1)!0.5!(oplus2)$);
        \newcommand{\oplusshortenTWO}{0pt};
        \node[oplussty, label={[xshift=0.5mm, yshift=0mm]center:$\binop_3$}]
                (oplus3) [right=of oplusmid] {}
                edge [pre, out=90, in=0, shorten >= \oplusshortenTWO] (oplus1) {}
                edge [pre, out=270, in=0, shorten >= \oplusshortenTWO] (oplus2) {};

        \node (ourstcqv) [right=of oplus3, xshift=-0.5cm] {} edge [pre] (oplus3) {};

        \pgfgetlastxy{\XCoord}{\YCoord};
        \outveccuboid{\XCoord + 5mm}{\YCoord}

        \newcommand{\leftlinexshift}{-2mm}
        \coordinate (left line bottom) at
                ($(current bounding box.south west) + (\leftlinexshift, 0)$);
        \coordinate (left line top) at
                ($(current bounding box.north west) + (\leftlinexshift, 0)$);
        \draw[semithick] (left line bottom) -- (left line top);

        \node[xshift=-3mm, yshift=-20mm, rotate=90]
                (mutan label) at (left line top) {MUTAN};

        \node[xshift=-3mm, yshift=15mm, rotate=90]
                (mutan label) at (left line bottom) {Ours};
\end{tikzpicture}

\end{center}

Our fusion operator generalizes MUTAN by allowing different
nonlinearities~$f_{rq}$ and~$f_{rv}$ to act on the features~$M_r\mutanfeat{q}$
and~$N_r\mutanfeat{v}$, followed with a neural network, producing a set of~$R$
output features, as in MUTAN after the Hadamard product step~$M_r\mutanfeat{q}
\odot N_r\mutanfeat{v}$.  In the case of MUTAN, the~$R$ output feature vectors
are combined by element-wise summation, whereas in our fusion operator the~$R$
feature vectors are combined by applying a tree of binary operators~$\binopb$.
}

\headerbox{5.~Ablation on Configurations}{name=ablation,span=1,column=0,below=tucker-decomp}{%
We compare performance improvements between different instantiations of the
generalized fusion operator, with the best one shown below.

\hspace{-2mm}\resizebox{1.05\linewidth}{!}{\newcommand{\branchstub}[5]{%
\begin{scope}[yshift=#2]
\node (q#1) {$M_{#5}\mutanfeat{q}$};
\node (v#1) [below=of q#1] {$N_{#5}\mutanfeat{v}$};

\node[hadamard, xshift=2cm] (h#1) at ($(q#1)!0.5!(v#1)$) {$\odot$}
         edge [pre] node [below] {$f_{\mathrm{selu}}$} (v#1) {}
         edge [pre] node [above, xshift=#4] {$f_{\mathrm{#3}}$} (q#1) {};

\node[nn] (nn#1) [right=of h#1] {$\Phi_#5$}
        edge [pre] (h#1) {};
\end{scope}}

\newcommand{\vecelems}[1]{%
        \node [matrix, vecsty]
                (vector#1) [right=of nn#1, xshift=1cm]
                {%
                        \node[rectangle] (vec#1a) {$t_1$};\\
                        \node[rectangle] (vec#1b) {$t_2$};\\[-2mm]
                        \node[rectangle] (vec#1c) [xshift=1mm] {$\vdots$};\\
                        \node[rectangle] (vec#1d) {$t_o$};\\
                };

        \draw [post, shorten >= 3pt, red!80, thick] (nn#1) -- (vector#1);
        \fill[color=red!65] ($(nn#1)!0.5!(vector#1)$) -- (vector#1.north west) -- (vector#1.south west) -- cycle;
}

\begin{tikzpicture}
        [gate/.style={circle,
                      inner sep=0pt,
                      minimum size=2mm,
                      thick,
                      draw=black,
                      fill=white},
         gateedge/.style={-{Circle[fill=white]},
                          out=0,
                          in=180,
                          shorten >= 3pt,
                          semithick},
         hadamard/.style={circle,
                          inner sep=0pt,
                          minimum size=6mm,
                          thick,
                          draw=blue!75,
                          fill=blue!20},
         nn/.style={rectangle,
                    draw=black!50,
                    fill=black!20,
                    thick,
                    inner sep=0pt,
                    minimum size=6mm},
         pre/.style={<-, shorten <= 1pt, >=stealth, semithick},
         post/.style={->, shorten >= 1pt, >=stealth, semithick},
         sigmoid/.style={circle,
                         inner sep=0pt,
                         minimum size=4mm,
                         thick,
                         draw=black,
                         fill=green!50,
                         draw=green!80},
         sumsty/.style={circle,
                        inner sep=0pt,
                        minimum size=6mm,
                        thick,
                        draw=orange!80,
                        fill=orange!50},
         vecsty/.style={fill=red!50, draw=red!80, very thick}]

        \branchstub{1}{0cm}{tanh}{1mm}{1}
        \branchstub{2}{-2.5cm}{identity}{2mm}{\sigma}

        \branchstub{3}{-5cm}{sigmoid}{2mm}{2}

        \vecelems{1}
        \vecelems{3}

        \node[sigmoid] (sigmoid2)
                [right=of nn2, xshift=-5mm] {$\sigma$}
                edge [pre] (nn2) {}
                edge [gateedge] (vec1a) {}
                edge [gateedge] (vec1b) {}
                edge [gateedge] (vec1d) {}
                edge [gateedge] (vec3a) {}
                edge [gateedge] (vec3b) {}
                edge [gateedge] (vec3d) {};

        \node[sumsty] (sum) [right=of sigmoid2, xshift=2cm] {$\Sigma$}
              edge [pre, out=90, in=0] (vector1) {}
              edge [pre, out=270, in=0] (vector3) {};

        \node (tcqv) [right=of sum] {$\T_c^{\q\v}$}
                edge [pre] (sum) {};

        \node[matrix] (legend) at (current bounding box.east) [xshift=3.5cm] {%
                \node[label=right:Question vector features.] {$\q$};\\[2mm]
                \node[label=right:Image vector features.] {$\v$};\\[2mm]
                \node[hadamard, label=right:Element-wise multiplication.] {$\odot$};\\[2mm]
                \node[nn, label=right:Neural network layers~${(\phi_l)}_r$.] {$\Phi_r$};\\[2mm]
                \node[sigmoid, label=right:Logistic sigmoid nonlinearity.] {$\sigma$};\\[2mm]
                \node[gate, label=right:Scalar multiplication.] {};\\[2mm]
                \node[vecsty, label=right: Elements of~$\tuckbranch \in \mathbb{R}^{t_o}$.] {$t_i$};\\[2mm]
                \node[sumsty, label=right: Element-wise sum.] {$\Sigma$};\\
        };

        \begin{scope}[on background layer]
                \node (r1) [fill=black!10, rounded corners, fit=(legend)] {};
        \end{scope}
\end{tikzpicture}
}

\begin{centering}
\begin{tabular}{rc}
\textbf{Model} & \textbf{VQA 1.0 val} \\
\midrule
MUTAN~\cite{ben2017mutan} & 61.54 \\
\midrule
Nonlinearity Ensembling (NE) & 61.66 \\
Feature Gating (FG) & 61.72 \\
NE + Polarity Swap (PS) & 61.77 \\
\textbf{NE + FG} & \textbf{61.86} \\
\end{tabular}
\end{centering}
}

\headerbox{6.~Results}{name=results,span=2,column=1,below=comparison-mutan}{%
We evaluate our best model on VQA~2.0 test-dev, and compare to previous state
of the art models upon which our work is based.

\begin{center}
% \begin{table*}[!t]
% \centering
% \caption{A comparison with the state of the art of our best single model on the
%          VQA~2.0 test-dev and test-std sets.}
\begin{tabular}{l*{5}{c}}
\centering
& \multicolumn{4}{c}{\textbf{VQA~2.0 test-dev}} \\
\textbf{Model} & All & Y/N & Number & Other \\
\midrule
MCB~\cite{DBLP:conf/emnlp/FukuiPYRDR16} & 61.96 & 78.41 & 38.81 & 53.23 \\
MUTAN~\cite{ben2017mutan} as trained and evaluated by us & 63.13 & 80.7 & 39.4 & 53.55 \\
ResNet features~$7 \times 7$ (single)~\cite{teney2017tips} & 62.07 & 79.20 & 39.46 & 52.62 \\
\midrule
Ours (single) & 64.22 & 81.19 & 40.95 & 55.05 \\
\midrule
\end{tabular}
% \label{tab:sota-comparison}
% \end{table*}
\end{center}

We attribute our improvement in accuracy to the introduction of nonlinearities
in the fusion operator.
}

\headerbox{7.~Conclusions}{name=conclusion,column=1,below=results,span=2,above=bottom}{%
\begin{itemize}
        \item Generalized the MUTAN multimodal fusion operator.

        \item We found configurations that demonstrated significant gains
                relative to MUTAN and MCB on the VQA 2.0 task.

        \item Our future work is focused on leveraging this generalization as a
                design space for architecture search by reinforcement learning
                or other methods.
\end{itemize}
}

\headerbox{8.~References}{name=references,column=0,span=1,below=ablation}{%

%\small % Reduce the font size in this block
\renewcommand{\section}[2]{\vskip 0.05em} % Get rid of the default "References" section title
%\nocite{*} % Insert publications even if they are not cited in the poster


\bibliographystyle{unsrt}
\bibliography{poster} % Use sample.bib as the bibliography file
}

\headerbox{9.~Acknowledgements}{name=acknowledgements,column=0,span=1,below=references,above=bottom}{


%\small % Reduce the font size in this block
\renewcommand{\section}[2]{\vskip 0.05em} % Get rid of the default "References" section title
%\nocite{*} % Insert publications even if they are not cited in the poster

This work received funding from CIFAR and NSERC\@.
}

\end{poster}

\end{document}
